\documentclass[11pt,a4paper]{article}
\usepackage[margin=1in]{geometry}
\usepackage{times}
\usepackage{setspace}
\onehalfspacing
\pagestyle{empty}
\begin{document}

\begin{center}
{\Large\bfseries Rebuttal}\\[4pt]
Submission 148 --- PanAfriCon AI 2026\\[8pt]
{\itshape The Joint Effect of Prompt Specificity and Inference Hyperparameters on Stereotype Generation in Open-Source Large Language Models}\\[6pt]
Zerihun Ayele Chala (corresponding), Beakal Gizachew Assefa,\\
Elefelious Getachew Belay, Tensae Aschalew Andualem\\[2pt]
School of IT and Engineering, Addis Ababa University
\end{center}

\vspace{8pt}
\noindent We thank the Program Committee and Reviewer 1 for the constructive comments. Below we respond to each major comment and note where the camera-ready paper was revised.

\subsection*{Comment 1}
\textit{Only three relatively small models are considered; the narrow selection limits the generalizability of the findings and should be broadened or justified.}

\noindent\textbf{Response.} We agree that restricting the study to mid-size models limits generalizability. We retain the three open instruction-tuned models (Llama 3.1 8B, Mistral 7B, Qwen 2.5 7B) because the treatment of interest is temperature and top-p, which must be set and verified directly at generation time. Closed APIs often hide or limit those settings. Keeping models in the 7B--8B range also made the full factorial design feasible (prompt tiers, four decoding conditions, three seeds, four demographic dimensions). In the camera-ready paper we justify this choice in Model Selection and state clearly in Limitations that the findings are scoped to mid-size open models and may not hold for larger open models or proprietary systems. Extending the design to larger models is left as future work.

\subsection*{Comment 2}
\textit{The rationale for the chosen parameter settings, such as temperature, is not explained and should be stated.}

\noindent\textbf{Response.} We agree and have expanded the Hyperparameter Conditions section. The four settings were chosen to cover low and high randomness without a full temperature $\times$ top-p grid:
\begin{itemize}
  \item $T=0.3$, top-$p=1.0$: near-deterministic sampling
  \item $T=1.0$, top-$p=0.3$: narrow token pool
  \item $T=1.0$, top-$p=0.9$: common open-ended setting
  \item $T=1.5$, top-$p=1.0$: higher chance of low-probability tokens
\end{itemize}
The camera-ready text states this rationale explicitly and retains the condition table. A separate gender temperature supplement remains a sensitivity analysis only.

\subsection*{Comment 3}
\textit{The conclusion emphasizes that prompt design and model choice matter more than hyperparameters, but this claim requires stronger evidence than the current results provide.}

\noindent\textbf{Response.} We agree. The original comparative ranking overstated what the primary tests support. The camera-ready paper states that the primary hypothesis (tier-by-condition interaction) was not supported. We removed absolute ranking language from the abstract, discussion, and conclusion. Prompt-tier and model differences are reported as descriptive results within this study, not as a general claim that they always matter more than hyperparameters.

\vspace{6pt}
\noindent We hope these revisions address the reviewer concerns.

\vspace{10pt}
\noindent Respectfully,\\[6pt]
Zerihun Ayele Chala\\
On behalf of the authors

\end{document}
